\documentclass[11pt]{article}

% Package Header

\usepackage[paperheight=11in,
			paperwidth=8.5in,
			inner=0.4in,
			textheight=8.9in,
			textwidth=7.8in,
			marginparwidth=120pt,
			marginparsep=32pt,
			bottom=1.4in,
			footskip=1.0in]{geometry}
\geometry{letterpaper}                   % ... or a4paper or a5paper or ...

\usepackage{array}
\usepackage{amssymb}
\usepackage{asymptote}
\usepackage{amsmath}
\usepackage{amsfonts}
\usepackage{animate}
\usepackage{arrayjobx}
\usepackage{arydshln}
\usepackage{bookman}
\usepackage{color}
%\usepackage{colortbl}
\usepackage{cancel}
\usepackage{calculator}
\usepackage{collcell}
\usepackage{epstopdf}
\usepackage{enumerate}
\usepackage{epsfig}
%\usepackage{etex}
\usepackage[english]{babel}
\usepackage{etoolbox}
% \usepackage{fancyhdr}
\usepackage{fancybox}
\usepackage[nomessages]{fp}
\usepackage{graphics}
\usepackage{graphicx}
\usepackage{ifthen}
\usepackage{ifmtarg}
\usepackage[latin1]{inputenc}
\usepackage{latexsym}
\usepackage{multirow}
\usepackage{multicol}
\usepackage{multimedia}
\RequirePackage[noplaybutton]{media9}
\usepackage{mathbbol}
\usepackage{mathtools}
\usepackage{multido}
\usepackage{palatino}
\usepackage{polynom}
\usepackage{pxfonts}
\usepackage{pifont}
\usepackage{pgffor}
%\usepackage{pgfplots}
\usepackage{pdfpages}
\usepackage[parfill]{parskip}
\usepackage{soul}
\usepackage{tikz}
\usetikzlibrary{arrows,shapes,trees,shapes.geometric,calc}
\usepackage{txfonts}
\usepackage{tcolorbox}
\usepackage{wasysym}
\usepackage{wrapfig}
\usepackage{xcolor}
\usepackage{xmpmulti}
\usepackage{xparse}
\usepackage{xstring}
\usepackage{enumitem}

%\usepackage{hyperref}
\usepackage{geometry}
\usepackage{amsmath}
\usepackage{amssymb}
\usepackage{amsthm}
\usepackage{amscd}
\usepackage{amsfonts}
\usepackage{graphicx}%
\usepackage{fancyhdr}
\usepackage{array}
\usepackage{pgf,tikz}
\usetikzlibrary{arrows}


\usepackage[osf]{mathpazo}

\geometry{textheight=10.0in, textwidth=7.7in}

\usepackage{sectsty}
\subsectionfont{\rmfamily\mdseries\itshape\large}

% Don't indent paragraphs.
\setlength\parindent{0em}


\begin{document}

\newcommand{\NP}{\newpage\thispagestyle{empty}}
\newcommand{\B}{\textcolor{blue}}
\newcommand{\R}{\textcolor{red}}
\newcommand{\G}{\textcolor{olive}}
\newcommand{\Or}{\textcolor{YellowOrange}}
\newcommand{\DEF}{\underline{Defn}:\ \ \ }
\newcommand{\THM}{\underline{Thm}:\ \ \ }
\newcommand{\PROP}{\underline{Properties}:\ \ \ }
\newcommand{\NOTE}{\underline{Notes}:\ \ \ }
\newcommand{\EX}{\underline{Ex}:\ \ \ }
\newcommand{\SOL}{\underline{Soln}:\ \ \ }
\newcommand{\UL}{\underline}
\newcommand{\RE}{\mathbf{R}}
\newcommand{\V}{\textbf{ v }}
\newcommand{\VV}{[v_1,v_2]}
\newcommand{\VVV}{[v_1,v_2,v_3]}
\newcommand{\W}{\textbf{ w }}
\newcommand{\WW}{[w_1,w_2]}
\newcommand{\WWW}{[w_1,w_2,w_3]}
\newcommand{\VA}{\textbf{ a }}
%\newcommand{\AA}{[a_1,a_2]}
\newcommand{\AAA}{[a_1,a_2,a_3]}
\newcommand{\VB}{\textbf{ b }}
\newcommand{\BB}{[b_1,b_2]}
\newcommand{\BBB}{[b_1,b_2,b_3]}
\newcommand{\VC}{\textbf{ c }}
\newcommand{\CC}{[c_1,c_2]}
\newcommand{\CCC}{[c_1,c_2,c_3]}
\newcommand{\VI}{\textbf{ i }}
\newcommand{\VJ}{\textbf{ j }}
\newcommand{\VK}{\textbf{ k }}
\newcommand{\VF}{\textbf{ F }}
\newcommand{\VG}{\textbf{ G }}
\newcommand{\VO}{{\vec{0}}}
\newcommand{\OBUL}{\quad $\circ$ \ \ \ }
\newcommand{\NEWLINE}{\\ \hspace{0.63cm}}
\newcommand{\CAN}{\cancel}
\newcommand{\GF}{{\nabla f}}
\newcommand{\PFPX}{{\dfrac{\partial f}{\partial x}}}
\newcommand{\PFPY}{{\dfrac{\partial f}{\partial y}}}
\newcommand{\PFPZ}{{\dfrac{\partial f}{\partial z}}}
\newcommand{\DFDX}{{\dfrac{df}{dx}}}
\newcommand{\GG}{{\nabla g}}
\newcommand{\PGPX}{{\dfrac{\partial g}{\partial x}}}
\newcommand{\PGPY}{{\dfrac{\partial g}{\partial y}}}
\newcommand{\PGPZ}{{\dfrac{\partial g}{\partial z}}}
\newcommand{\DGDX}{{\dfrac{dg}{dx}}}
\newcommand{\PPX}{{\dfrac{\partial}{\partial x}}}
\newcommand{\PPY}{{\dfrac{\partial}{\partial y}}}
\newcommand{\PPZ}{{\dfrac{\partial}{\partial z}}}
\newcommand{\PP}{\partial}
\newcommand{\DELF}{\nabla \cdot \textbf{ F }}
\newcommand{\CURLF}{\nabla \times \textbf{ F }}
\newcommand{\DEL}{\nabla \cdot }
\newcommand{\CURL}{\nabla \times }
\newcommand{\RT}{\textbf{ r }(t)}
\newcommand{\RPT}{\textbf{ r }'(t)}
\newcommand{\GINT}{\iint \limits_{R} (\text{curl} \ \textbf{ F }) \cdot \textbf{ k } \ dA}
\newcommand{\ON}{\onslide}
\newcommand{\OB}{\overbrace}
\newcommand{\OBT}{\overbracket}
\newcommand{\OS}{\overset}
\newcommand{\UB}{\underbrace}
\newcommand{\UBT}{\underbracket}
\newcommand{\US}{\underset}
\newcommand{\D}{\displaystyle}
\newcommand{\MCol}{\multicolumn}
\newcommand{\RRE}{\mathbf{R}^2}
\newcommand{\RRRE}{\mathbf{R}^3}
\newcommand{\DXDT}{{\dfrac{dx}{dt}}}
\newcommand{\DYDT}{{\dfrac{dy}{dt}}}
\newcommand{\DZDT}{{\dfrac{dz}{dt}}}
\newcommand{\FRT}{\textbf{ F }(\textbf{ r }(t))}
\newcommand{\WORK}{ {\int \limits_{C} \textbf{ F } \cdot d\textbf{ r }} }
\newcommand{\WORKAB}{ {\int \limits_{A}^{B}\textbf{ F } \cdot d\textbf{ r }} }
\newcommand{\WORKT}{ {\int \limits_{a}^{b}\textbf{ F }(\textbf{ r }(t)) \cdot \textbf{ r }'(t) \ dt} }
\newcommand{\FLUX}{ {\iint \limits_{S}\textbf{ F } \cdot d\textbf{ S }} }
\newcommand{\FLUXX}{ {\iint \limits_{S}\textbf{ F } \cdot \textbf{ n } \ dS} }
\newcommand{\FLUXXX}{ {\iint \limits_{R}\textbf{ F }(\textbf{ r }(u,v)) \cdot \textbf{ N } \ du\ dv} }
\newcommand{\OWORK}{ {\oint \limits_{C}\textbf{ F } \cdot d\textbf{ r }} }
\newcommand{\GREENT}{ {\oint \limits_{C}\textbf{ F }\cdot \textbf{ T } \ ds} }
\newcommand{\GREENN}{ {\oint \limits_{C}\textbf{ F }\cdot \textbf{ N } \ ds} }
\newcommand{\DIVTHM}{\iiint \limits_{T} \text{div} \ \textbf{ F }  \ dV}
\newcommand{\GREENCURL}{\iint \limits_{R} (\text{curl} \ \textbf{ F }) \cdot \textbf{ k } \ dA}
\newcommand{\STOKECURL}{\iint \limits_{S} (\text{curl} \ \textbf{ F }) \cdot \textbf{ n } \ dS}
\newcommand{\GREENDIV}{\iint \limits_{R} \text{div} \ \textbf{ F } \ dA}
\newcommand{\GREENSA}{\iint \limits_{R} \left(\dfrac{\partial N}{\partial x} - \dfrac{\partial M}{\partial y}\right) \ dA}
\newcommand{\GREENSL}{\oint \limits_{C} (M\ dx + N\ dy )}
\newcommand{\FDR}{ {\textbf{ F } \cdot d\textbf{ r }} }
\newcommand{\FDRT}{ {\textbf{ F }(\textbf{ r }(t)) \cdot \textbf{ r }'(t)} }
\newcommand{\BCOLS}{\begin{columns}}
\newcommand{\BCOL}{\begin{column}}
\newcommand{\ECOLS}{\end{columns}}
\newcommand{\ECOL}{\end{column}}
\newcommand{\BLINE}{\tikz[baseline]}
\newcommand{\NODERRB}{\node[fill=red!10,draw,rectangle,anchor=base]}
\newcommand{\NODEGRB}{\node[fill=green!10,draw,rectangle,anchor=base]}
\newcommand{\NODEBRB}{\node[fill=blue!10,draw,rectangle,anchor=base]}
\newcommand{\NODECRB}{\node[fill=red!0,draw,rectangle,anchor=base]}
\newcommand{\NODERR}{\node[fill=red!10,rectangle,anchor=base]}
\newcommand{\NODEGR}{\node[fill=green!10,rectangle,anchor=base]}
\newcommand{\NODEBR}{\node[fill=blue!10,rectangle,anchor=base]}
\newcommand{\NODECR}{\node[fill=red!0,rectangle,anchor=base]}
\newcommand{\NODERE}{\node[fill=red!10,draw,ellipse,anchor=base]}
\newcommand{\NODEGE}{\node[fill=green!10,draw,ellipse,anchor=base]}
\newcommand{\NODEBE}{\node[fill=blue!10,draw,ellipse,anchor=base]}
\newcommand{\NODECE}{\node[fill=red!0,draw,ellipse,anchor=base]}
\newcommand{\NODERC}{\node[fill=red!10,draw,circle,anchor=base]}
\newcommand{\NODEGC}{\node[fill=green!10,draw,circle,anchor=base]}
\newcommand{\NODEBC}{\node[fill=blue!10,draw,circle,anchor=base]}
\newcommand{\NODECC}{\node[fill=blue!0,draw,circle,anchor=base]}
\newcommand{\CNODE}{\tikz[na]\node [coordinate] }
\newcommand{\BENDL}{edge [bend left] }
\newcommand{\BENDR}{edge [bend right] }
\newcommand{\ARROWR}{\path [red,->]}
\newcommand{\ARROWG}{\path [green,->]}
\newcommand{\ARROWB}{\path [blue,->]}
\newcommand{\ARROWSR}{\path [red,<->]}
\newcommand{\ARROWSG}{\path [green,<->]}
\newcommand{\ARROWSB}{\path [blue,<->]}
\newcommand{\LINER}{\path [red,-]}
\newcommand{\LINEG}{\path [green,--]}
\newcommand{\LINEB}{\path [blue,--]}
\newcommand{\DE}{east}
\newcommand{\DW}{west}
\newcommand{\DN}{north}
\newcommand{\DS}{south}

\newcommand{\TT}[1]{\texttt{#1}}

\tikzstyle{every picture}+=[remember picture]
\tikzstyle{na} = [baseline=-1.5ex]

%%-----------------------------------------------------------
%% Bliss Definitions
\newcommand{\To}{\Rightarrow}
\newcommand{\del}{\nabla}
\newcommand{\C}{\mathcal{C}}
\newcommand{\N}{\mathcal{N}}
\newcommand{\A}{\mathcal{A}}
\renewcommand{\O}{\mathcal{O}}
% \renewcommand{\N}{\mathbb{N}}
\newcommand{\ben}{\begin{enumerate}}
\newcommand{\een}{\end{enumerate}}
\newcommand{\eps}{\varepsilon}
\renewcommand{\epsilon}{\varepsilon}
\newcommand{\p}{\partial}
\newcommand{\<}{\langle}
\renewcommand{\>}{\rangle}
\renewcommand{\subset}{\subseteq}
\newcommand{\cont}{\Rightarrow \Leftarrow}
\newcommand{\back}{\backslash}
\newcommand{\abs}[1]{|{#1}|}
\newcommand{\ip}[1]{\langle{#1}\rangle}
\newcommand{\m}[1]{{\mathbf{#1}}}
\newcommand{\sts}{\m{S}^\textrm{T}\m{S}}
\newcommand{\rank}{\textrm{rank}}
\newcommand{\X}{\mathcal{X}}
\newcommand{\Y}{\mathcal{Y}}
\def\ds{\displaystyle}
\def\d{\partial}
\renewcommand{\arraystretch}{1.5}
%%-----------------------------------------------------------


\newcommand\CEN[1]{\multicolumn{1}{c}{ #1 }}

\newcommand\MATL[3]
{{%
\arraycolsep=#1\arraycolsep\ensuremath{{\setlength{\extrarowheight}{#2cm}\begin{bmatrix*}[l]#3\end{bmatrix*}}}
}}
\newcommand\MATC[3]
{{%
\arraycolsep=#1\arraycolsep\ensuremath{{\setlength{\extrarowheight}{#2cm}\begin{bmatrix*}[c]#3\end{bmatrix*}}}
}}
\newcommand\MATR[3]
{{%
\arraycolsep=#1\arraycolsep\ensuremath{{\setlength{\extrarowheight}{#2cm}\begin{bmatrix*}[r]#3\end{bmatrix*}}}
}}

\newcommand\MATDL[3]
{{%
\arraycolsep=#1\arraycolsep\ensuremath{{\setlength{\extrarowheight}{#2cm}\begin{vmatrix*}[l]#3\end{vmatrix*}}}
}}
\newcommand\MATDC[3]
{{%
\arraycolsep=#1\arraycolsep\ensuremath{{\setlength{\extrarowheight}{#2cm}\begin{vmatrix*}[c]#3\end{vmatrix*}}}
}}
\newcommand\MATDR[3]
{{%
\arraycolsep=#1\arraycolsep\ensuremath{{\setlength{\extrarowheight}{#2cm}\begin{vmatrix*}[r]#3\end{vmatrix*}}}
}}

\newcommand\TABwB[3]
{{
\setlength{\extrarowheight}{#1 cm}
\begin{tabular}{#2}
	\hline
	#3
	\hline
\end{tabular}
}}

\newcommand\TABwoB[3]
{{
\setlength{\extrarowheight}{#1cm}
\begin{tabular}{#2}
	#3
\end{tabular}
}}

\newcommand\includegraphicsDD[2]
{%
\includegraphics[scale=#1]{#2}
}

\newcommand\includegraphicsDDD[2]
{%
\includemedia[scale=#1,3Dmenu,activate=onclick,deactivate=onclick,
3Droll=0.,
3Dortho=0.0046491301618516445,
3Dc2c=0.8617081642150879 -0.29301872849464417 0.4142450988292694,
3Dcoo=-13.165555000305176 3.489872455596924 52.97535705566406,
3Droo=149.99999574342704,
3Dlights=Headlamp,add3Djscript=asylabels.js]{\includegraphics{#2}}
{#2.prc}
}

\newcommand\MATS[2]
{{%
  \arraycolsep=#1\arraycolsep\ensuremath{\begin{bmatrix*}[r]#2\end{bmatrix*}}
}}

\newcommand{\OKFACE}{
\begin{tikzpicture}
\draw[thick] (0cm,0cm) circle(.2cm);
\draw[thick] (-0.1,-0.09) -- (0.1,-0.09);
\draw [thick, fill=black] (0,0) circle (0.01);
\draw [rotate=90, fill=white] (0.06,0.06) circle (0.03);
\draw [rotate=90, fill=white] (0.06,-0.06) circle (0.03);
\end{tikzpicture}
}

\newcommand{\HAPPYFACE}{
\begin{tikzpicture}
\draw[thick] (0cm,0cm) circle(.2cm);
\draw[thick] plot [smooth,tension=1.5] coordinates{(-0.1,-0.08) (0,-0.13) (0.1,-0.08)};
\draw [thick, fill=black] (0,0) circle (0.01);
\draw [rotate=90, fill=white] (0.06,0.06) circle (0.04);
\draw [rotate=90, fill=white] (0.06,-0.06) circle (0.04);
\end{tikzpicture}
}

\newcommand{\SADFACE}{
\begin{tikzpicture}
\draw[thick] (0cm,0cm) circle(.2cm);
\draw[thick] plot [smooth,tension=1.5] coordinates{(-0.1,-0.12) (0,-0.07) (0.1,-0.12)};
\draw [thick, fill=black] (0,0) circle (0.01);
\draw [rotate=90, fill=white] (0.06,0.06) circle (0.02);
\draw [rotate=90, fill=white] (0.06,-0.06) circle (0.02);
\end{tikzpicture}
}

\newcommand{\CounterReset}{
\setcounter{CritCounter}{1}
\setcounter{TestCounter}{1}
\setcounter{ColorCounter}{1}
\setcounter{Index}{1}
}

\newcommand\CritLineInitialize[1]{
\ifdef{\SymFun}
{
	\setcounter{CritLineSlide}{#1}
}
{
	\newarray\TestPts
	\newarray\TestVals
	\newarray\TestEVal
	\newarray\fIncDec
	\newarray\testsigns
	\newarray\inequality
	\setcounter{CritLineSlide}{#1}
	\newcommand\SymFun{}
	\newcommand\IncDecStr{}
	\newcommand\TestSignStr{}
	\newcommand\InEqStr{}
	\newcommand\SymFunD{}
	\newcommand\SymFunDLt{}
	\newcommand\SymFunN{}
	\newcommand\SymFunNLt{}
	\newcommand{\CritLineTable}{}
}
}

\newcommand\USubInitialize{
	\newcommand{\Fun}{}
	\newcommand{\FunN}{}
	\newcommand{\FunD}{}
	\newcommand{\FunU}{}
	\newcommand{\UIntegrand}{}
	\newcommand{\AFunU}{}
	\newcommand{\DX}{}
}

\newcommand*\circled[1]{\tikz[baseline=(char.base)]{
            \node[shape=circle,draw,inner sep=1pt,color=blue] (char) {#1};}}

\newcommand*\ellipsed[1]{\tikz[baseline=(char.base)]{
            \node[shape=ellipse,draw,inner sep=1pt,color=blue] (char) {#1};}}

\newcommand*\rected[1]{\tikz[baseline=(char.base)]{
            \node[shape=rectangle,draw,inner sep=1pt,color=blue] (char) {#1};}}

\newcommand*\starred[1]{\tikz[baseline=(char.base)]{
            \node[shape=star,star points=10,draw,inner sep=1pt,color=blue] (char) {#1};}}

\newcommand\CHECKBOX[1]{\makebox[0pt][l]{$\square$}\raisebox{.15ex}{\hspace{0.1em}\ON#1{\B{$\checkmark$}}}}

\makeatletter
\newcommand{\@RemovePositiveSignIfAny}[1]{\IfBeginWith{#1}{+}{\StrGobbleLeft{#1}{1}}{#1}}%
\newcommand{\@PlusIfNonZero}[1]{\IfStrEq{#1}{}{}{+}}%
\newlength{\@MySystemWidth}%
\newcommand*{\@ChooseSignWithAlignment}[2]{%
    \IfBeginWith{#2}{-}%
    {{}-\makebox[\@MySystemWidth][#1]{\ensuremath{\StrGobbleLeft{#2}{1}}}}%
    {{}\@PlusIfNonZero{#2}\makebox[\@MySystemWidth][#1]{\ensuremath{\@RemovePositiveSignIfAny{#2}}}}%
}%
\newcommand*{\@ChooseSignWithRightAlignment}[1]{\@ChooseSignWithAlignment{r}{#1}}%
\newcommand*{\@ChooseSignWithLeftAlignment}[1]{\@ChooseSignWithAlignment{l}{#1}}%

\newenvironment{MySystem}[3][1.2em]{% [<width>] {num of variables} {alignment}
    \newcolumntype{R}{>{\collectcell\@ChooseSignWithRightAlignment}{r}<{\endcollectcell}}
    \newcolumntype{L}{>{\collectcell\@ChooseSignWithLeftAlignment}{l}<{\endcollectcell}}
    \setlength{\@MySystemWidth}{#1}% Store the width to use to typeset the terms
    \IfStrEqCase{#3}{%
        {R}{\begin{array}{r@{} *{\numexpr#2-1\relax}{#3@{}} @{{}={}} r}}%
        {L}{\begin{array}{l@{} *{\numexpr#2-1\relax}{#3@{}} @{{}={}} l}}%
    }[\parbox{\linewidth}{\color{red}MySystem Error: Only L and R alignment types are supported. (using R)}\begin{array}{r *{\numexpr#2-1\relax}{#3} @{{}={}} r}]%
}{%
    \end{array}%
}
\newenvironment{MySystemNoEqual}[3][1.2em]{% [<width>] {num of variables} {alignment}
    \newcolumntype{R}{>{\collectcell\@ChooseSignWithRightAlignment}{r}<{\endcollectcell}}
    \newcolumntype{L}{>{\collectcell\@ChooseSignWithLeftAlignment}{l}<{\endcollectcell}}
    \setlength{\@MySystemWidth}{#1}% Store the width to use to typeset the terms
    \IfStrEqCase{#3}{%
        {R}{\begin{array}{r@{} *{\numexpr#2\relax}{#3@{}} @{{}{}} r}}%
        {L}{\begin{array}{l@{} *{\numexpr#2\relax}{#3@{}} @{{}{}} l}}%
    }[\parbox{\linewidth}{\color{red}MySystem Error: Only L and R alignment types are supported. (using R)}\begin{array}{r *{\numexpr#2\relax}{#3} @{{}{}} r}]%
}{%
    \end{array}%
}

\makeatother


$\quad$

\vspace{2.6cm}

\begin{center}
\LARGE {MA 301} \\

\Large{Test 2} \\

\large{Spring - 2015} \\

\vspace{1.3cm}

\setlength{\extrarowheight}{0.4cm}
\begin{tabular}{|c|c|c|}
\hline
Problem & points & out of  \\
\hline
\hline
1 & & 6 \\
\hline
2 & & 5 \\
\hline
3 & & 10 \\
\hline
4 & & 9 \\
\hline
\hline
Score & & 30\\
\hline
\end{tabular}

\vspace{1.5cm}

%--------TABLE---------
	\setlength{\extrarowheight}{0.1cm}
	\begin{tabular}{rcc}
		Name: & \underline{$\qquad \qquad \qquad \qquad \quad$}& \\
		&&$\hspace{1.4cm}$\\
		HR: & \underline{$\qquad \qquad \qquad \qquad \quad$}& \\
		&&\\
		Signature: & \underline{$\qquad \qquad \qquad \qquad \quad$}& \\
	\end{tabular}
%--------------------

\vspace{1.5cm}

\vspace{0.5cm}

	\setlength{\extrarowheight}{0.1cm}
	\begin{tabular}{| p{12cm} |}
		\hline
		\multicolumn{1}{|c|}{{\bf READ BEFORE YOU START}} \\
		\hline
		 * Show all work.\\
		 * Change your calculator mode to RADIANS.\\
		 * Sloppy solutions may be penalized.\\
		 * Use calculators ONLY for computations.\\
		 * You must write up every problem in a clear and \\
		 $\,\,$ organized way.\\
		 * If you write solutions on a page different than the \\
		 $\,\,$ question, then clearly state where it is located.\\
		\hline
	\end{tabular}

\end{center}

\vspace{1.0cm}

\thispagestyle{empty} 
\newpage
\thispagestyle{empty}

{\bf Problem 1:} Suppose {\bf F} $\D{= \left[\cos(\pi x),\, xz,\, \frac{1}{2}x^2+\ln y\right]}$ and $S$ the surface: $\D{z = 4-y^2}$ where $y \ge 0$ and $0\le x \le 3$ with outward pointing normal (away from the $x$-axis).  Let $C$ be the curve which bounds $S$.  

\vspace{0.2cm}

(a) \,\, Compute \, the line integral of ${\bf F}$ on $C$: $\D \OWORK$

\vspace{18.2cm}

(b) \,\, \UL{Set-up} the double integral that computes the flux of ${\bf F}$ through $S$: \, $\D \iint \limits_{S} {\bf F} \cdot {\bf n} \, dS$

\thispagestyle{empty}
\newpage
\thispagestyle{empty}

{\bf Problem 2:} Let $\text{\bf F} = [107,ze^{xz},xyz]$ and $S$ be the surface (with outward orientation) of the solid bounded by the surfaces $x = 0$, $x = 5$, $y = 0$, $z = 0$, and $y^2 + z^2 = 1$ in the first octant.

\vspace{0.2cm}

Compute \,$\D{\iint \limits_{S} \text{\bf F}\cdot \text{\bf n} \, dS}$

\thispagestyle{empty}
\newpage
\thispagestyle{empty}

%{\bf Problem 3:} The the closed cube-like surface, $S$, shown below has a top that is an unknown smooth surface. Let ${\bf F} = (x+2yz){\bf i} - 8y{\bf j} + (11z+\cos(x)){\bf k}$. Suppose we knew the outward flux through sides $A$ and $B$ were both $-4$ and that the volume of the solid is 10.  Find the outward flux through the top?	
%
%\begin{center}
%	\includegraphics[scale = 0.2]{../../../Figures/Test2_Fig1.png}
%\end{center}

{\bf Problem 3:} Let \,\,$z_1 = -3i,\,\, z_2 = 3-3i,\,\, z_3 = -5,\,\, z_4 = -\sqrt{3}-i$.  \quad ({\bf Do NOT leave any answers as decimals.}) 

\vspace{0.01cm}

{\it **Bonus: 2 additional bonus points will be given if you do all work using Euler's Polar Form (i.e. $z = re^{\theta i}$) on (b) - (d). }

\vspace{0.01cm}
 
(a) Perform the following operations and leave answer in \UL{standard form}.

\vspace{0.3cm}

\qquad $\triangleright \quad z_1 - z_4$

\vspace{1.0cm}

\qquad $\triangleright \quad \D{\frac{z_4}{z_2}}$

\vspace{1.0cm}

(b) Write $z_1, \, z_2, \, z_3, \, z_4$ in \UL{polar form}.

\vspace{0.3cm}

\qquad $\triangleright \quad z_1 = $

\vspace{1.2cm}

\qquad $\triangleright \quad z_2 = $

\vspace{1.2cm}

\qquad $\triangleright \quad z_3 = $

\vspace{1.2cm}

\qquad $\triangleright \quad z_4 = $

\vspace{1.2cm}

(c) Perform the following operations and leave answer in \UL{polar form}. 

\vspace{0.5cm}

\qquad $\triangleright \quad z_1z_2$

\vspace{1.0cm}

\qquad $\triangleright \quad (z_1z_2)^6$

\vspace{1.0cm}

(d) Solve the following complex equations for $x$ and leave your solutions in polar form: 

\qquad $\triangleright \quad \D{\frac{x^{1/6}}{z_1} - z_2 = 0}$.

\vspace{2.0cm}

\qquad $\triangleright \quad \D{x^3-z_3=0}$.

\thispagestyle{empty}
\newpage
\thispagestyle{empty}

{\bf Problem 4:}  Answer the following quick computation questions:

(a) \,\, Let $\D{{\bf F} = -y\,\VI + x\, \VJ + e^z\, \VK}$ and $S$ be a circle of radius 3 inside the plane $x+2y+z=10$.  If $C$ is the boundary of this circle (oriented counterclockwise) then find $\D{\OWORK}$ \,\,(your answer should be in terms of $a,b, \& \, c$)

\vspace{8.0cm}

(b) \,\, Compute $\D{\OWORK}$ where ${\bf F} = [5y ,2z ,5x]$ and $C$ is the boundary around the triangle obtained by tracing out the path $(2,0,0)$ to $(2,0,2)$, to $(2,2,2)$ back to $(2,0,0)$. 

\vspace{7.8cm}

(c) \,\, Compute $\D{\FLUXX}$ where ${\bf F} = [5x^2+y ,-3xy-3y,-6xz]$ and $S$ is the surface of the solid cube with $\D{0 \le x \le 2,}$ \, $\D{ |y| \le 1,}$ \, and $\D{|z| \le 1}$ with outward pointing normal.

%\thispagestyle{empty}
%\newpage
%\thispagestyle{empty}

%\noindent {\bf Problem 1:} Let \,\,$z_1 = -3i,\,\, z_2 = 3-3i,\,\, z_3 = -5,\,\, z_4 = -\sqrt{3}-i$.
%
%\noindent(a) Perform the following operations and leave answer in \UL{standard form}. ({\bf No Decimals})
%
%\vspace{0.3cm}
%
%\qquad $\triangleright \quad z_1 - z_4$
%
%\vspace{1.0cm}
%
%\qquad $\triangleright \quad \D{\frac{z_4}{z_2}}$
%
%\vspace{1.0cm}
%
%\noindent(b) Write $z_1, \, z_2, \, z_3, \, z_4$ in \UL{Euler's polar form}. ({\bf No Decimals})
%
%\vspace{0.3cm}
%
%\qquad $\triangleright \quad z_1 = $
%
%\vspace{1.2cm}
%
%\qquad $\triangleright \quad z_2 = $
%
%\vspace{1.2cm}
%
%\qquad $\triangleright \quad z_3 = $
%
%\vspace{1.2cm}
%
%\qquad $\triangleright \quad z_4 = $
%
%\vspace{1.2cm}
%
%\noindent(c) Perform the following operations and leave answer in \UL{Euler's polar form}. ({\bf No Decimals})
%
%\vspace{0.8cm}
%
%\qquad $\triangleright \quad z_1z_2$
%
%\vspace{1.5cm}
%
%\qquad $\triangleright \quad (z_1z_2)^6$
%
%\vspace{1.5cm}
%
%\noindent(d) Solve the complex equations for $x$ and leave your solutions in Euler's Form: $re^{\theta i}$. ({\bf No Decimals})
%
%\vspace{0.6cm}
%
%\qquad $\triangleright \quad \D{\frac{x^{1/6}}{z_1} - z_2 = 0}$.
%
%\vspace{2.5cm}
%
%\qquad $\triangleright \quad \D{x^3-z_3=0}$.
%
%\thispagestyle{empty}
%\newpage
%\thispagestyle{empty}
%
%\noindent{\bf Problem 2:} Let ${\bf F} = [-3xz^2, 6, 5x^2y]$ and $S$ is the parabolic cylinder $S:\, y=x^2$ where $0 \le x \le 2$ and $0 \le z \le 3$ and normal vector point away from the $y$-axis.  Compute the flux of ${\bf F}$ through the parabolic cylinder $S$
%
%\vspace{0.2cm}
%
%\thispagestyle{empty}
%\newpage
%\thispagestyle{empty}
%
%\noindent{\bf Problem 3:} Let ${\bf F}$ and $S$ be the same as in {\bf Problem 2}. If $C$ be the boundary of $S$.  Compute \quad $\D{\OWORK}$
%
%\thispagestyle{empty}
%\newpage
%\thispagestyle{empty}
%
%\noindent{\bf Problem 4:} Let $\text{\bf F} = [x,ye^z,e^z]$ and $S$ the surface (outward orientation) of the box $|x|\le 1$, $|y|\le 2$ , $-1\le z \le 7$. Compute the flux through $S$.
%
%\vspace{0.2cm}
%
%\thispagestyle{empty}
%\newpage
%\thispagestyle{empty}
%
%\noindent{\bf Problem 5:} Let $F = [x^3-y^3,y^3-z^3,z^3-x^3]$ and $S$ is the surface of the solid bounded by: 
%
%\vspace{0.2cm}
%
%%--------TABLE---------
%\setlength{\extrarowheight}{0.2cm}
%\begin{tabular}{cl}
%	(i) & quarter sphere $x^2 + y^2 + z^2 = 25$ \, ($y\ge0$, $z\ge 0$); \\
%
%	(ii) & the plane \, $z = 0$; \\
%
%	(iii) & the plane $y=0$ \\
%\end{tabular}
%%--------------------
%
%\vspace{0.4cm}
%
%\noindent with outward pointing normal. Compute the flux through $S$.
%
%\thispagestyle{empty}
%\newpage
%\thispagestyle{empty}
%
%{\bf Problem 3:} The the closed cube-like surface, $S$, shown below has a top that is an unknown smooth surface. Let ${\bf F} = (x+2yz){\bf i} - 8y{\bf j} + (11z+\cos(x)){\bf k}$. Suppose we knew the outward flux through sides $A$ and $B$ were both $-4$ and that the volume of the solid is 10.  Find the outward flux through the top?	
%
%\begin{center}
%	\includegraphics[scale = 0.2]{../../../Figures/Test2_Fig1.png}
%\end{center}
%
%\thispagestyle{empty}
%\newpage
%\thispagestyle{empty}
%
%\noindent{\bf Problem 6:}  Let $S$ be a surface given by $z=x^3y$ where $0 \le x \le 5$ and $-1 \le y \le 6$. {\bf Set-up} (DO NOT Compute) the double integral ({\it with limits of integration}) you would use to find the surface integral \quad $\D{\iint \limits_{S} (1 + 9xz)^{3/2}\, dS}$
%
%\thispagestyle{empty}
%\newpage
%\thispagestyle{empty}
%	
%Use Stoke's Theorem to evaluate $\displaystyle \int_{C} \mathbf{F} \cdot d\mathbf{r} where \mathbf{F}(x, y, z) = x\mathbf{i} + y\mathbf{j} + 1\!(x^{2} + y^{2})\mathbf{k} $ and $C$ is the boundary of the part of the paraboloid where $z = 4 - x^{2} - y^{2}$ which lies above the $xy$-plane and $C$ is oriented counterclockwise when viewed from above. 
%
%\hrule 
%
%Use Stokes' Theorem to find the line integral of ${\bf F} = [ xy, yz, xz ]$ around the boundary of the surface S given by $z = 25 - x^2$ for $0 \leq x \leq 5 and -5 \leq y \leq 5$, oriented upward. Sketch both S and its boundary C.
%
%$\displaystyle \int_C {\bf F} \cdot d{\bf r} =$
%
%\hrule 
%
%Use Stokes' Theorem to find the line integral ${\bf F} = 2 y {\bf i} + 7 z {\bf j} + 6 x {\bf k}$ around the triangle obtained by tracing out the path $(6,0,0)$ to $(6,0,2)$, to $(6,3,2)$ back to $(6,0,0)$.
%
%$\displaystyle \int_C {\bf F} \cdot d{\bf r} = $
%
%\hrule 
%
%Use Stokes' Theorem to find the line integral ${\bf F} = 2 y {\bf i} + 7 z {\bf j} + 6 x {\bf k}$ around a circle C of radius 7 centered at (9,3,1) in the plane $z = 1$, oriented counterclockwise when viewed from above.
%
%$\displaystyle \int_C {\bf F} \cdot d{\bf r} =$
%
%\hrule 
%
%Evaluate the line integral ${\bf G} = xy {\bf i} + z{\bf j} + 2 y{\bf k}$ around a square of side 4, centered at the origin, lying in the yz-plane, and oriented counterclockwise when viewed from the positive x-axis.
%
%$\displaystyle \int_C {\bf F} \cdot d{\bf r} =$
%
%\hrule 
%
%Find $\int_C{\bf F} \cdot d{\bf r}$ where C is a circle of radius 2 in the plane $x+y+z=2$, centered at (2,4,-4) and oriented clockwise when viewed from the origin, if \,\, ${\bf F} = -6 z{\bf j} + 7 y{\bf k}$
%
%$\int_C{\bf F} \cdot d{\bf r} =$
%
%\hrule 
%
%Use Stoke's theorem to evaluate $\displaystyle \int\!\!\int_{S} \text{curl}(\mathbf{F}) \cdot d\mathbf{S}$ where $\mathbf{F}(x, y, z) = -18\!yz\mathbf{i} + 18\!xz\mathbf{j} + 5\!(x^{2} + y^{2})z\mathbf{k}$ and S is the part of the paraboloid $z = x^{2} + y^{2}$ that lies inside the cylinder $x^{2} + y^{2} = 1$, oriented upward. 
%
%\thispagestyle{empty}
%\newpage
%\thispagestyle{empty}
%
%Evaluate $\displaystyle \iint \limits_S \mathbf{F} \cdot d \mathbf{A} where \mathbf{F} = (3xy^2, 3x^2y, z^3)$ and S is the surface of the sphere of radius 6 centered at the origin. 
%
%\hrule 
%
%Find the outward flux of the vector field $\mathbf{F} = [x^3, y^3, z^2]$ across the surface of the region that is enclosed by the circular cylinder $x^2 + y^2 = 49$ and the planes $z = 0$ and $z = 8$. 
%
%\hrule 
%
%Use the divergence theorem to calculate the flux of the vector field ${\bf F}(x,y,z) = 2 x y {\bf i} + z^{2} {\bf j} + 8 y {\bf k}$ through the outward-oriented, closed box S given by
%
%$0 \leq x \leq 5, \quad 0 \leq y \leq 2, \quad 0 \leq z \leq 2$.
%
%$\displaystyle \iint\limits_S {\bf F} \cdot d{\bf A} =$
%
%\thispagestyle{empty}
%\newpage
%\thispagestyle{empty}
%
% Let $f(x,y,z)=e^{-z}$ and the surface, S, be given by
%$x^2+y^2=16 \text{,} \quad 0\le z\le 9$
%
%$\displaystyle{\iint \limits_{\mathcal{S}} f(x,y,z)\,dS=} $
%
%\hrule 
%
% Let $f(x,y,z)=z$ and the surface, S, be given by
%$y = 1-z^2 \text{,} \qquad 0\le x \le 7 \quad \& \quad 0 \le z\le 7$
%
%$\displaystyle{\iint \limits_{\mathcal{S}} f(x,y,z)\,dS= }$
%
%\hrule 
%
%$\mathbf{F}(x,y,z) = [0,0,\cos(x^2+y^2)]$ and the surface, S, is the upward oriented disk given by
%
%$x^2+y^2 \leq 16, \quad z = 3$
%
%Flux of ${\bf F}$ through S =
%
%\hrule 
%
%Find the flux of the constant vector field $\mathbf{F} = -3\,\mathbf{i}-4\,\mathbf{j}+3\,\mathbf{k}$ through a square plate of area 25 in the $yx$-plane oriented in the positive z-direction.
%
%Flux =
%
%\hrule 
%
%The flux of the constant vector field a $\mathbf{i}+b\mathbf{j}+c\mathbf{k}$ through the square of side 7 in the plane y=6, oriented in the positive y-direction, is -196. Which of the constants a, b, c can be determined from the information given? Give the value(s).
%
%(If the value for a given constant is not determined, enter undetermined for your answer, below.)
%
%a =
%
%b =
%
%c =
%
%\hrule 
%
%Let $\mathbf{F} = 5 x\mathbf{i} + 5 y\mathbf{j}$ and S be the upward oriented surface given by
%
%$z=25 - (x^2+y^2), \quad z \ge 0$
%
%Flux of ${\bf F}$ through S = 
%
%\hrule 
%
%Let $\mathbf{F} = x \mathbf{i} + y \mathbf{j} + 3 z\mathbf{k}$ and S be the outward oriented quarter cylinder given by
%
%$x^2+y^2=16, \, 0\leq x\leq 4, \, 0\leq y\leq 4, \, 0\leq z\leq 4$
%
%Flux of ${\bf F}$ through S =
%
%\hrule 
%
%Find the surface area of the part of the sphere $x^{2} + y^{2} + z^{2} = 100$ that lies above the cone $z = \sqrt{x^{2} + y^{2}}$
%
%\hrule 
%
%If a parametric surface given by $\mathbf{r_{1}}(u, v) = f(u, v)\mathbf{i} + g(u, v)\mathbf{j} + h(u, v)\mathbf{k}$ and $-5 \leq u \leq 5, -3 \leq v \leq 3$, has surface area equal to 1, what is the surface area of the parametric surface given by $\mathbf{r_{2}}(u, v) = 1\mathbf{r_{1}}(u, v)$ with $-5 \leq u \leq 5, -3 \leq v \leq 3$? 
%
%\thispagestyle{empty}
%\newpage
%\thispagestyle{empty}
%

	
\end{document}